% !TEX root = chapter-standalone.tex



\section{Products}
\label{sec:arithmetic-products}


\subsection{Commutativity}

\todotext{JL: rewrite this section}

The equation $x \cdot y = y \cdot x$ describes the \emph{commutativity property} of multiplication.

Let us look at the corresponding relationship on the level of sets:
\begin{equation}\label{eq:mult-commutativity-sets}
    \setA \cartprod \setB \simeq \setB \cartprod \setA.
\end{equation}
The canonical isomorphism in this case, which gives proof for this relationship, is
\begin{equation}\label{eq:braiding-sets}
    \defmapperiod{
        \braiding
    }{
        \setA \cartprod \setB
    }{
        \to
    }{
        \setB \cartprod \setA
    }{
        \tup{\ela, \elb}
    }{
        \tup{\elb, \ela}
    }
\end{equation}
We use the same notation again, $\braiding$, and also call this isomorphism a ``braiding''.

This function is ``canonical'' because, once again, it's ``shape'' or ``recipe'' works for any two sets $\setA$ and $\setB$.
To explicate this in more detail, consider, for example, the particular sets $\setA = \makeset{0, 1, 2}$ and $\setB = \makeset{\star, \dagger}$.
Then, besides the canonical function $\braiding \colon \setA \cartprod \setB \to \setB \cartprod \setA$, there are also \emph{other} (non-canonical) isomorphisms possible, such as for instance the function $\setA \cartprod \setB \to \setB \cartprod \setA$ which maps like this:
\begin{align}\label{eq:noncanonical-braiding}
    \tup{0, \star} \mapsto \tup{\star, 1} \\
    \tup{1, \star} \mapsto \tup{\star, 2} \\
    \tup{2, \star} \mapsto \tup{\star, 3} \\
    \tup{0, \dagger} \mapsto \tup{\dagger, 1} \\
    \tup{1, \dagger} \mapsto \tup{\dagger, 2} \\
    \tup{2, \dagger} \mapsto \tup{\dagger, 3}.
\end{align}
It may be written more compactly by the recipe ``$\tup{x, y} \mapsto \tup{y, x +1 \text{ mod } 3}$".
However, this recipe still depends on specific features of the set $\setA$ and would not work for any two arbitrary sets (namely, it uses the fact that the elements of $\setA$ in this example are integers for which the operation ``mod 3'' makes sense).
Thus this function is not ``canonical'' (while $\braiding$ is).


\subsection{Associativity}

\todotext{JL: rewrite this section}

The equations $x \cdot (y \cdot z) = (x \cdot y) \cdot z$ describes the \emph{associativity} property of multiplication of natural numbers.
Correspondingly, there is a canonical isomorphism
\begin{equation}\label{eq:sets-prod-associator}
    \associator \colon (\setA \cartprod \setB) \cartprod \setC \to \setA \cartprod (\setB \cartprod \setC)
\end{equation}

This isomorphism $\associator$ is
\begin{equation}\label{eq:sets-prod-associator-def}
    \defmapperiod{
        \associator
    }{
        (\setA \cartprod \setB) \cartprod \setC
    }{
        \to
    }{
        \setA \cartprod (\setB \cartprod \setC)
    }{
        \tup{\tup{\ela, \elb}, \elc}
    }{
        \tup{\ela, \tup{\elb, \elc}}
    }
\end{equation}


\todotext{J: add remark about how the two three-way cartesian products defined by binary operation relate to the trinary 'balanced' cartesian product}

\subsection{Singletons are like the number one}

\todotext{@JL: write up}

In \cref{sec:special-sets} we defined the singleton set $\singleton = \makeset{\singletonel}$ in order to have -- for practical reasons of exposition and readability -- a ``standard, go-to, default one-element set''.
This set (and any other singleton set) behaves like the natural number ``$1$'' in the sense that
\begin{equation}\label{eq:singleton-set-is-neutral-for-prod}
    \singleton \cartprod \setA \isomorphic \setA \quad \text{and} \quad \setA \cartprod \singleton \isomorphic \setA
\end{equation}
is true for any set $\setA$. In each case there is also a \emph{canonical} isomorphism, one which does not involve any ad hoc choices:
\begin{equation}\label{eq:leftunitor-sets-prod}
    \defmapcomma{
        \leftunitor
    }{
        \singleton \cartprod \setA
    }{
        \to
    }{
        \setA
    }{
        \tup{\singletonel, \ela}
    }{
        \ela
    }
\end{equation}
and
\begin{equation}\label{eq:rightunitor-sets-prod}
    \defmapperiod{
        \rightunitor
    }{
        \setA \cartprod \singleton
    }{
        \to
    }{
        \setA
    }{
        \tup{\ela, \singletonel}
    }{
        \ela
    }
\end{equation}
The names $\leftunitor$ and $\rightunitor$ stand for \emph{left unitor} and \emph{right unitor}, respectively.
These functions are canonical because their ``shape'' or ``recipe'' does not depend on the particular set $\setA$ involved; the recipe is simply ``forget about the element $\singletonel$".
It is in this sense that it works for ``all cases at once'' and does not depend on any extra or ad hoc choices.


\subsection{The empty set is (again) like the number zero}

The relationship
\begin{equation}\label{eq:emptyset-is-zero}
    \Emptyset \cartprod \setA  = \Emptyset,
\end{equation}
which holds for any set $\setA$, is analogous to the equation
\begin{equation}
    0 \cdot x = 0
\end{equation}
for any natural number $x$. (And also $\setA \cartprod \Emptyset  = \Emptyset$, just as $x \cdot 0 = 0$.)

These equations follow from \cref{def:cartesian-product} (or from the formulation given in \cref{rem:formal-cartesian}).

\todotext{J: make a remark why this equation is put together with distributivity...}

That this relationship is an equality -- and not merely an isomorphism like all the other ``arithmetic relationships'' between sets that we will address -- has to do with the fact that there exists only one unique set having the cardinality zero, namely the empty, while for any non-zero natural number $x$, there exist infinitely-many sets having the cardinality $x$.



\subsection{From products to exponents}

\todotext{JL: this subsection still needs its proper write-up; what follows is only a first bridging paragraph together with the simplest of the exponent relationships, stated as an exercise.}

In \cref{sec:functions} we introduced the notation $\setB^{\setA}$ for the set of all functions from $\setA$ to $\setB$, and we motivated the exponent notation by the observation that, for finite sets, $\cardof{\setB^{\setA}} = {\cardof \setB}^{\cardof \setA}$.
The arithmetic analogy goes further than a statement about sizes: the familiar rules for exponents have counterparts which hold for arbitrary sets, in the form of isomorphisms.
The simplest of these is the counterpart of the rule $x^1 = x$.
\begin{exercise}
    \label{ex:functions-out-of-a-singleton}
    Let $\setB$ be any set, and let $\singleton = \makeset{\singletonel}$ be the singleton set.
    Show that
    \begin{equation}
        \label{eq:exponent-singleton}
        \setB^{\singleton} \isomorphic \setB,
    \end{equation}
    by exhibiting a function $\setB^{\singleton} \sto \setB$ together with an inverse to it (\cref{def:function-isomorphism}).
\end{exercise}

\begin{solution}
    We define a function
    \begin{equation}
        \defmapcomma{
            \stylemorph{\Phi}
        }{
            \setB^{\singleton}
        }{
            \sto
        }{
            \setB
        }{
            \mapa
        }{
            \mapa(\singletonel)
        }
    \end{equation}
    which evaluates a function at the unique element of $\singleton$, and a function
    \begin{equation}
        \defmapcomma{
            \stylemorph{\Psi}
        }{
            \setB
        }{
            \sto
        }{
            \setB^{\singleton}
        }{
            \elb
        }{
            \mapa_{\elb}
        }
    \end{equation}
    where $\mapa_{\elb} \colon \singleton \sto \setB$ is defined by $\mapa_{\elb}(\singletonel) \definedas \elb$.
    This last prescription really does define a function: $\singletonel$ is the only element of $\singleton$, so a value has been specified at every element of the source, and unambiguously so.

    It remains to verify the two equations of \cref{def:function-isomorphism}.
    For any $\elb \setin \setB$ we have
    \begin{equation}
        (\stylemorph{\Psi} \mthen \stylemorph{\Phi})(\elb) = \stylemorph{\Phi}(\mapa_{\elb}) = \mapa_{\elb}(\singletonel) = \elb,
    \end{equation}
    so $\stylemorph{\Psi} \mthen \stylemorph{\Phi} = \catidat{\setB}$.
    Conversely, for any function $\mapa \colon \singleton \sto \setB$ we have
    \begin{equation}
        (\stylemorph{\Phi} \mthen \stylemorph{\Psi})(\mapa) = \stylemorph{\Psi}(\mapa(\singletonel)) = \mapa_{\mapa(\singletonel)},
    \end{equation}
    and this function agrees with $\mapa$ at $\singletonel$, hence at every element of $\singleton$.
    Two functions with the same source and target which agree at every element of the source are equal: by \cref{def:function} they consist of exactly the same ordered pairs.
    Therefore $\stylemorph{\Phi} \mthen \stylemorph{\Psi} = \catidat{\setB^{\singleton}}$, and $\stylemorph{\Phi}$ is an isomorphism.
\end{solution}

Note that \cref{eq:exponent-singleton} makes precise the observation, made in \cref{sec:functions-working}, that specifying a function out of a singleton set amounts to nothing more than choosing a single element of the target.
Note also the difference in role played by the singleton set here and in \cref{eq:singleton-set-is-neutral-for-prod}: there it was a \emph{factor} that could be discarded, while here it is an \emph{exponent}, that is, an index set with exactly one index.

\subsection{Arbitrary dependent products}

\todotext{Rewrite below to make clear the distinction between cartesian product as defined via tuples and the dependent product defined via functions}

Given any finite number of sets $\setAn{1}, \dots, \setAn{n}$, we have defined their cartesian product as
\begin{equation}
    \setAn{1} \cartprod \setAn{2} \cartprod \dots \cartprod \setAn{n} = \makeset{ \tup{\ela_{1}, \ela_{2}, \dots, \ela_{n}} \mid \ela_{i} \setin \setAn{i} \text{ for each } i = 1, \dots, n }.
\end{equation}
This defines an ``ordered'' cartesian product because tuples arrange elements in an ordered sequence: the \emph{first} entry is associated with the index ``1'', the \emph{second} entry with the index ``2'', and so on, all the way up to $n$.
With this definition, the cartesian product $\setA \cartprod \setB \cartprod \setC$ is not, for example, equal to the cartesian product $\setC \cartprod \setB \cartprod \setA$.
The order of the factors matters.

There is, however, an unordered version of the cartesian product.
To explain how it works, we start with the observation that we can think of a tuple
\begin{equation}
    \tup{\ela_{1}, \ela_{2}, \dots, \ela_{n}}
\end{equation}
(with $\ela_{i} \setin \setAn{i}$ for each $i = 1, \dots, n$) as encoding (and as encoded by) a function
\begin{equation}
    \mapa : \makeset{1, \dots, n} \mto \bigsetunion_{\setIel \setin\setI } \setA_{\setIel}
\end{equation}
such that $\mapa(\setIel) \setin\setA_{\setIel}$ for each $\setIel \setin \makeset{1, \dots, n}$.
Namely, $\mapa(\setIel) = \ela_\setIel$, for each $\setIel \setin \makeset{1, \dots, n}$.
In this description of the tuple via a function, the ``ordering'' is encoded via the fact that the elements of $\makeset{1, \dots, n}$ are ordinal numbers.

We can however, also consider any family of sets $\makeset{ \setA_\setIel }_{\setIel \setin\setI}$ for an \emph{arbitrary} set $\setI$ which serves as an index set (previously we were only using the index set $\setI = \makeset{1, \dots, n}$).
In this case, we define the unordered cartesian product of the family of sets $\makeset{ \setA_\setIel }_{\setIel \setin\setI}$ by
\begin{equation}
    \underset{\setIel \setin\setI}{\stylesets{\prod}} \setA_{\setIel} \definedas \makeset{ \mapa \colon \setI \to \bigsetunion_{\setIel \setin\setI } \setA_{\setIel} \mid \mapa(\setIel) \setin\setA_{\setIel} \ \forall \setIel \setin\setI}.
\end{equation}
We use the symbol~$\stylesets{\prod}$ instead of the infix notation~``$\cartprod$'' (which necessarily requires an ordering).
The upper case greek letter pi $\prod$ is meant to stand for ``product''.
Note that this unordered definition does not make any assumptions on the size of $\setI$; it can also be infinite.




